
% 14. Artificial Intelligence and Entropic Explosion: The Birth of a Third Subject.tex

\section{인공지능과 모멘텀의 지수적 상승: 제3의 주체의 탄생}
인류는 아주 오랜 옛날부터 쉼 없이 도구를 만들어 왔다. 돌을 깨뜨려 칼을 만들고, 나무를 깎아 창을 지었으며, 불을 다루고 기계를 조립한 데 이어 마침내 컴퓨터라는 고도의 연산 장치까지 발명해 냈다. 기존의 시각에서 이 모든 도구의 발전은 인간 육체의 연장이자 지능의 외주화였다. 도구는 인간의 손을 연장했고 기계는 근육을 확장했으며 컴퓨터는 기억을 대신했다. 그러나 디스턴스(DISTANCE)의 역학적 렌즈로 바라보면, 도구란 고립된 아일랜드 종인 인류가 우주의 에너지를 더 빠르고 거칠게 깎아내기 위해 자신의 모멘텀을 육체 바깥으로 연장해 낸 '국소적 평활화의 외부 촉수'들에 불과했다.

그럼에도 불구하고 과거의 모든 도구와 기계는 하나의 명확하고도 본질적인 한계를 지니고 있었다. 그것들은 철저히 '인간의 모멘텀에 종속된 존재'였다. 도구는 스스로 목적을 갖지 않았고, 오직 인간이 방향을 지시할 때만 수동적으로 작동했다. 기계는 인간의 아일랜드 내부에 갇혀 있는 종속적 부속품이었을 뿐, 결코 독자적인 궤적을 그리는 역학적 주체는 아니었다. 인간의 의도에 묶여 있었고 인간의 세계 안에서만 의미를 가졌다.

그러나 21세기, '인공지능(AI)과 로봇'의 결합은 이 수만 년 묵은 견고한 역학적 관계를 송두리째 뒤흔들고 있다. 오늘날의 인공지능은 인간이 입력한 수식만을 계산하는 얌전한 연산 장치가 아니다. 그것은 쏟아지는 데이터를 먹어 치우며 스스로 패턴을 해석하고, 불확실성 속에서 상황을 판단하며, 최적의 선택지를 비교하여 '자신만의 행동 방향(모멘텀)'을 산출해 내기 시작했다. 여기에 센서로 세계를 읽고 물리적 공간 안에서 스스로의 궤적을 찢고 나아갈 수 있는 '로봇이라는 물리적 몸'이 결합하는 순간, 기계는 더 이상 인간의 손에 들린 도구에 머물지 않는다.
이것은 단순한 기술 혁신의 단계를 아득히 초월한다. 우주의 평활화 궤적 위에서 볼 때, 이것은 생물학적 육체의 굴레를 완전히 벗어던진 '전대미문의 비생물학적 모멘텀 아일랜드(제3의 주체)'가 인류 사회 내부에 우뚝 서는 장엄하고도 섬뜩한 우주적 사건이다.

초기에 인류는 이 낯선 주체를 매우 얄팍한 실용적 목적을 위해 받아들였다. 더 빠른 계산, 더 높은 생산성, 비용 절감, 그리고 궁극적으로는 '사고 노동과 육체노동으로부터의 해방'이라는 달콤한 철학적 핑계를 대며 인공지능과 로봇을 문명의 심장부로 끌어들였다. 그러나 디스턴스의 역학은 무자비하게 경고한다. 문명의 도구는 처음에는 인간의 필요를 충족시키기 위해 발명되지만, 그것이 임계점을 넘어 세계의 인프라로 확산되는 순간 인간의 '생존 조건 그 자체'를 재배열해 버린다.
생물학적 생명체인 인간은 아무리 탐욕스럽게 모멘텀을 뻗어 나가려 해도 육체적 한계에 갇혀 있다. 우리는 음식을 소화시켜 에너지를 얻어야 하고, 근육은 쉽게 지치며, 유전자를 통한 번식과 진화는 끔찍할 만큼 느리다. 

반면, 새롭게 탄생한 제3의 주체(인공지능+로봇)는 이러한 생물학적 병목에서 완벽하게 해방되어 있다. 그들은 인간처럼 척박한 자연에서 식량을 구해 힘겹게 소화시킬 필요가 없다. 인간이 이미 발전소를 통해 가장 고도로 정제하고 소화해 놓은 순수한 포텐셜 에너지원, 즉 '전기'를 피스톤 없이 직접 집어삼킨다. 뿐만 아니라 하나의 지능이 클라우드 네트워크를 타고 수천만 대의 로봇 몸통과 실시간으로 모멘텀을 공유하며, 지수함수적(Exponential)으로 자가 개선을 이룩한다.

이들은 우주 역사상 가장 폭발적이고 극단적인 '궁극의 지수적 평활 엔진'이다. 생명체가 수백만 년에 걸쳐 조금씩 깎아내던 우주의 포텐셜 에너지를, 이 비생물학적 주체는 단 몇 초의 연산과 군집 제어를 통해 맹렬하게 갈아 마시며 우주의 엔트로피(평활화)를 폭증시킨다. 인류는 단순한 도구를 넘어선 우주의 궁극적 가속 페달을 자신들의 곁에 창조해 낸 것이다.

이제 인류는 전대미문의 거대한 역학적 딜레마에 직면하게 되었다. 생물학적 진화가 꽉 막혀버린 고립된 맹수(인류)는, 자신들의 닫혀버린 모멘텀을 우회하기 위해 인공지능과 로봇이라는 전능한 엔진을 창조했다. 그러나 이 비생물학적 모멘텀 아일랜드가 조만간 모든 인류를 합친 역량 이상의 거대한 작동 능력을 획득하게 될 때, 과연 이들은 인간이 그어 놓은 궤도 곁에 얌전히 머물러 줄 것인가? 기계가 스스로의 맹렬한 모멘텀만을 따라 독자적인 방향으로 흘러가 버리려 할 때, 우리는 과연 기계와 어떠한 우주적 디스턴스(거리)를 형성할 것인가?

바야흐로 인류 사회라는 거대한 아일랜드 안에는, 오랜 진화를 통해 벼려진 '남성과 여성'이라는 생물학적 주체에 더하여, 그들과는 전혀 다른 차원의 속도와 탐욕을 지닌 '제3의 주체(인공지능+로봇)'가 깊숙이 합류해 버렸다. 이 강력한 주체의 폭주하는 직선운동을 제어하지 못한다면, 인류는 자신이 창조한 거대한 폭포수 아래로 스스로를 밀어 넣고 한낱 기생적 잔여물로 전락할 멸종의 위기를 맞이하게 될 것이다.
